\documentclass[journal]{IEEEtran}
\usepackage{graphicx}
\usepackage{amsmath}
\usepackage{amssymb}
\usepackage{booktabs}
\usepackage{multirow}
\usepackage{url}
\usepackage{cite}
\usepackage{xcolor}

% ---------------------------------------------------------------
% TODO markers: search for "TODO" throughout this file.
% Every TODO marks content that MUST be filled from your real
% results once training on real/annotated data is complete.
% Do NOT submit with any TODO still in the text.
% ---------------------------------------------------------------

\begin{document}

\title{Multi-Class Hate Speech Detection in Sinhala-English Code-Mixed
Social Media Text with Emoji Feature Integration Using XLM-RoBERTa}

\author{H.D. Nethmini Hakmanage, A. Senuri Kodithuwakku
\thanks{H.D. Nethmini Hakmanage and A. Senuri Kodithuwakku are with the
Faculty of Information Technology, Horizon Campus, Sri Lanka.}}

\markboth{IEEE Access}{Hakmanage \MakeLowercase{\textit{et al.}}: Multi-Class Hate Speech Detection in Sinhala-English Code-Mixed Text}

\maketitle

\begin{abstract}
Hate speech detection in low-resource, code-mixed languages remains an
open challenge, particularly for Sinhala-English social media text where
prior work has been limited to binary classification, has discarded
emoji information during preprocessing, and has been restricted to a
single platform. We address these three gaps by fine-tuning
XLM-RoBERTa-base for three-class classification (hate, offensive,
neutral) on Sinhala-English code-mixed comments collected from Facebook
and YouTube, comparing a baseline model against a proposed model that
encodes emoji as textual sentiment cues. % TODO: extend abstract with
% one-sentence summary of your actual quantitative result once available,
% e.g. "The proposed model achieved a macro-F1 of X.XX (+/- X.XX),
% compared to X.XX for the baseline (Wilcoxon p = X.XX)."
We describe our data collection, annotation, and evaluation methodology
in detail, including a rigorous data-quality audit process, and report
our findings with an explicit discussion of the failure modes and
limitations encountered.
\end{abstract}

\begin{IEEEkeywords}
Hate speech detection, code-mixed text, Sinhala-English, XLM-RoBERTa,
emoji features, low-resource NLP, social media.
\end{IEEEkeywords}

\section{Introduction}

With the growing popularity of social media, hate speech detection in
low-resource languages has emerged as an urgent research problem. In Sri
Lanka, Facebook and YouTube provide large volumes of user-generated,
code-mixed Sinhala text written in Roman characters interspersed with
English. Muthuthanthri and Smith \cite{muthuthanthri2024} built an
initial baseline dataset of 2,205 labelled Sinhala-English Facebook
comments and evaluated nineteen machine learning algorithms for hate
versus non-hate classification, with BERT achieving 82\% accuracy and a
90\% ROC-AUC. Liyanage and Jayakumar \cite{liyanage2021} similarly
observed a performance ceiling for classical machine learning models on
Romanised Sinhala text, and Hettiarachchi et al. \cite{hettiarachchi2020}
confirmed that pre-trained embeddings outperform bag-of-words
representations for this task.

Muthuthanthri and Smith \cite{muthuthanthri2024} explicitly identified
three limitations for future work: integrating emoji and symbol
information into detection models, developing methods to counter
symbol-based evasion, and extending data collection beyond Facebook to
other platforms such as YouTube. To our knowledge, no subsequent study
of Sinhala-English code-mixed hate speech has addressed these three gaps
together. This paper makes the following contributions:

\begin{itemize}
\item We extend the classification scheme from binary (hate / not-hate)
to a three-class taxonomy (hate, offensive, neutral) following Davidson
et al. \cite{davidson2017}, distinguishing targeted hate speech from
general offensive language that is not directed at an identity group.
\item We integrate emoji information as an explicit input feature via
Unicode-to-text description encoding, rather than discarding it during
preprocessing as in \cite{muthuthanthri2024}.
\item We collect data across two platforms (Facebook and YouTube)
rather than a single platform, following the cross-platform
generalisation argument supported by Vasantharajan and Thayasivam
\cite{vasantharajan2022}.
\end{itemize}

Our research question is: \textit{does fine-tuning XLM-RoBERTa with
emoji encoding lead to a statistically significant improvement in
macro-averaged F1-score over a fine-tuned XLM-RoBERTa baseline without
emoji encoding, on three-class Sinhala-English code-mixed data from
Facebook and YouTube, evaluated across five stratified folds?} This
question is falsifiable: if emoji features introduce noise, cause
overfitting on the minority class, or fail to produce a significant
macro-F1 improvement, the hypothesis is disproved.

The remainder of this paper is organised as follows. Section II reviews
related work in Sinhala and code-mixed hate speech detection and
multilingual Transformer models. Section III describes our data
collection, annotation, preprocessing, and model architecture in detail.
Section IV reports our results across both experimental conditions.
Section V discusses these results, including a dedicated failure
analysis and ethical reflection. Section VI concludes with our
contributions in light of the results and directions for future work.

% TODO: This section may need a closing paragraph summarising your
% actual finding once real results are available, to preview the paper's
% conclusion for the reader (standard IMRaD practice).

\section{Related Work}

Early classical machine learning approaches to Sinhala and Singlish hate
speech detection established meaningful but limited baselines.
Hettiarachchi et al. \cite{hettiarachchi2020} and Liyanage and Jayakumar
\cite{liyanage2021} both found that ensemble classifiers over
hand-crafted features plateau around 84\% accuracy on Romanised Sinhala
text, with no transformer-based methods explored in either study.
Fernando et al. \cite{fernando2022} showed that FastText combined with
RNNs achieves strong AUC-ROC on Sinhala social media text, motivating the
move toward learned embeddings over classical features.

Muthuthanthri and Smith \cite{muthuthanthri2024} remain the closest prior
work to ours: a Facebook-only, binary-labelled dataset evaluated with
BERT, explicitly excluding emoji from the input. Davidson et al.
\cite{davidson2017} demonstrated, in the (monolingual) hate speech
literature more broadly, that collapsing hate and offensive speech into
a single category measurably reduces classifier accuracy, motivating our
adoption of their three-class taxonomy. Sap et al. \cite{sap2019} showed
that hate speech classifiers trained on majority-dialect data
systematically misclassify minority-dialect text, which directly informs
our ethical considerations regarding potential bias against Tamil and
Muslim community speech in the Sri Lankan context.

On the modelling side, Conneau et al. \cite{conneau2020} introduced
XLM-RoBERTa, pretrained on 2.5TB of multilingual Common Crawl data across
100 languages, and Vasantharajan and Thayasivam \cite{vasantharajan2022}
demonstrated its superiority over BERT and SVM baselines on the
linguistically similar task of Tamil-English code-mixed YouTube comment
classification -- directly motivating both our choice of XLM-RoBERTa and
our inclusion of YouTube as a data source. Mozafari et al.
\cite{mozafari2020} further showed that models pretrained on web-crawled
text outperform those trained on curated corpora for social media
classification tasks, reinforcing the choice of XLM-RoBERTa (Common
Crawl-pretrained) over standard BERT. Barbieri et al. \cite{barbieri2020}
established macro-F1 as the standard metric for imbalanced multi-class
tweet classification benchmarks, which we adopt as our primary metric.

Across this body of work, the most advanced prior Sinhala-English
system -- Muthuthanthri and Smith \cite{muthuthanthri2024} -- remains
constrained to binary classification, discards emoji during
preprocessing, and draws from a single platform. No subsequent study
combines a three-class taxonomy, emoji-aware input encoding, and
cross-platform data collection for this language pair; this is the
specific, compound gap our work addresses.

\section{Methodology}

\subsection{Data Collection}

Data was collected from two social media platforms: Facebook, via the
Facepager tool interfacing with the Facebook Graph API, and YouTube, via
the YouTube Data API v3. % TODO: replace this sentence with an accurate
% description of what was ACTUALLY collected for your final results --
% e.g. if Facebook collection was not completed due to Meta's Graph API
% access restrictions (a documented, legitimate constraint -- see
% Section~\ref{sec:limitations}), state that plainly here rather than
% implying both platforms were used if only YouTube data was used.
Collection targeted public comment threads on political, religious, and
news-related content between % TODO: insert your actual collection date
% range %, expected to yield high-volume code-mixed Sinhala-English
commentary. Only publicly available posts and comments were collected;
private groups, individual pages, and direct messages were excluded.

\subsection{Annotation}
% TODO: replace this paragraph with your ACTUAL annotation process. If
% you used 2 annotators with an agreement-based (rather than 3-annotator
% majority-vote) protocol due to time constraints, say so explicitly --
% this is a legitimate, documented methodological choice, not something
% to obscure.

Each collected comment was independently labelled following a written
annotation guideline defining three classes -- hate speech, offensive
language, and neutral content -- with worked examples drawn from prior
literature \cite{muthuthanthri2024,davidson2017}. % TODO: state your
% actual number of annotators and reconciliation method here, and report
% your actual inter-annotator agreement statistic (Cohen's Kappa for 3+
% annotators, or percent agreement / Cohen's Kappa for 2 annotators).

\subsection{Preprocessing}

Comments were preprocessed via URL removal, username/mention removal,
emoji extraction and encoding (or removal, for the baseline condition),
lowercasing of English tokens, repeated-character normalisation, and
whitespace normalisation, with a minimum length filter of three tokens
applied after preprocessing. Emoji were converted to their Unicode text
description (e.g. an angry-face emoji becomes the text ``angry face'')
using the
Python \texttt{emoji} library and appended inline to the comment text,
allowing the model to learn emoji meaning jointly with surrounding
context rather than treating it as noise.

\subsection{Data Quality Audit}
\label{sec:audit}

Prior to model training, the annotated corpus was audited for exact and
near-duplicate content using core-text clustering (filler-word and
emoji-stripped normalisation followed by exact and TF-IDF-based
near-duplicate grouping). % TODO: report your actual final dataset's
% audit statistics here -- number of rows, number of distinct core
% clusters, duplicate/near-duplicate rate. This audit step is a genuine
% methodological strength if your final dataset passes it cleanly;
% report it honestly either way.
This step is critical for any code-mixed social-media hate speech
corpus, particularly given the risk that near-duplicate or templated
content can trivially inflate held-out performance metrics through
train/test leakage rather than genuine model generalisation -- a failure
mode we discuss in detail in Section~\ref{sec:failure-analysis}.

\subsection{Model Architecture}

The system fine-tunes XLM-RoBERTa-base \cite{conneau2020} for three-class
sequence classification. The only architectural difference between the
baseline and proposed conditions is the emoji-handling step in
preprocessing (Section III-C): the baseline strips emoji entirely,
matching the approach of \cite{muthuthanthri2024}, while the proposed
model encodes emoji as appended text tokens. All other architectural and
training choices are held constant across both conditions so that any
observed difference in macro-F1 can be attributed specifically to emoji
handling. Table~\ref{tab:training-config} summarises the training
configuration.

\begin{table}[h]
\centering
\caption{Training Configuration}
\label{tab:training-config}
\begin{tabular}{@{}ll@{}}
\toprule
Parameter & Value \\
\midrule
Base model & xlm-roberta-base \\
Max sequence length & 128 tokens \\
Optimiser & AdamW \\
Learning rate & 2e-5 \\
Dropout & 0.1 \\
Class weighting & Inverse-frequency balanced \\
Early stopping & Validation-loss based, patience 2 \\
\bottomrule
\end{tabular}
\end{table}

Class-weighted cross-entropy loss was used to address label imbalance,
with weights computed as the inverse frequency of each class in the
training fold. AdamW was selected over standard SGD because it decouples
weight decay from the gradient update, which has been shown to improve
generalisation when fine-tuning large pretrained Transformer models
\cite{conneau2020}; a learning rate of 2e-5 follows the range
established as effective for fine-tuning BERT-family models on
downstream classification tasks. XLM-RoBERTa was chosen over a
convolutional or recurrent architecture because the task requires
transferring pretrained multilingual representations to a low-resource,
code-mixed setting with limited labelled data -- a setting where
Transformer-based transfer learning has been shown to substantially
outperform architectures trained from scratch on code-mixed South Asian
language data \cite{vasantharajan2022,mozafari2020}.

\subsection{System Architecture}

Figure~\ref{fig:architecture} shows the end-to-end pipeline: raw
platform data is collected, annotated by independent human raters,
passed through the preprocessing pipeline described above, tokenised,
and fed to the fine-tuned XLM-RoBERTa classifier, which outputs a
three-class softmax distribution over \{neutral, offensive, hate\}.

\begin{figure}[h]
\centering
% TODO: insert your actual system architecture diagram image here, e.g.:
% \includegraphics[width=0.9\linewidth]{system_architecture.png}
\fbox{\parbox{0.85\linewidth}{\centering\vspace{1em}
Facebook + YouTube $\rightarrow$ Raw Corpus $\rightarrow$ Annotation +
Preprocessing $\rightarrow$ Feature Extraction $\rightarrow$
XLM-RoBERTa-base $\rightarrow$ Output Classification $\rightarrow$
Evaluation\vspace{1em}}}
\caption{End-to-end system architecture for multi-class Sinhala-English
code-mixed hate speech detection with emoji feature integration.}
\label{fig:architecture}
\end{figure}

\subsection{Baseline Definition}

The baseline is an XLM-RoBERTa-base model trained identically to the
proposed model in every respect except that emoji are stripped from the
input rather than encoded. This isolates emoji handling as the sole
experimental variable, allowing any difference in macro-F1 between the
two conditions to be attributed specifically to emoji information.

\subsection{Evaluation Protocol}

Macro-averaged F1-score is the primary evaluation metric, selected
because it weights all three classes equally regardless of frequency --
appropriate given the class imbalance in hate speech corpora, where
accuracy is a misleading metric \cite{muthuthanthri2024}, and consistent
with its use as the standard metric for imbalanced multi-class tweet
classification \cite{barbieri2020}. Secondary metrics include per-class
precision, recall, and F1, macro-averaged AUC-ROC, and confusion
matrices for failure analysis. Evaluation uses stratified $k$-fold cross-validation ($k=5$) rather than
a single fixed train/validation/test split, because with a target corpus
of approximately 2,500 comments a single split would leave only
approximately 375 comments for testing -- too few to yield a stable
per-class F1 estimate given the three-way class imbalance. Five-fold
cross-validation instead produces five independent test estimates,
allowing performance to be reported as a mean and standard deviation
that reflects genuine variability rather than the outcome of one
arbitrary split, with stratification ensuring the class distribution is
preserved in every fold. % TODO: state your actual k-fold value
% (5-fold as originally planned, or fewer if adjusted for a smaller
% real dataset) and confirm stratification / leakage-safety approach.
The identical fold assignment is reused for both the baseline and
proposed conditions to enable a paired statistical comparison. Statistical significance between
the two conditions' per-fold macro-F1 scores is assessed using the
Wilcoxon signed-rank test, chosen over a paired $t$-test because it does
not assume normally distributed differences, which is appropriate given
the small number of paired observations.

\section{Results}
\label{sec:results}

% TODO: THIS ENTIRE SECTION MUST BE REPLACED WITH YOUR REAL RESULTS.
% Do not submit with placeholder text below. Structure to follow:
%   - Table of mean +/- std macro-F1 for baseline vs proposed, per fold
%   - Per-class precision/recall/F1 table
%   - Confusion matrices (both conditions)
%   - Wilcoxon test statistic and p-value, with the n=5 (or n=k) low
%     statistical power caveat stated explicitly
% Example table structure (replace values):

\begin{table}[h]
\centering
\caption{Macro-F1 by Condition (TODO: replace with real values)}
\label{tab:results}
\begin{tabular}{@{}lcc@{}}
\toprule
Condition & Mean Macro-F1 & Std. Dev. \\
\midrule
Baseline (no emoji) & TODO & TODO \\
Proposed (with emoji) & TODO & TODO \\
\bottomrule
\end{tabular}
\end{table}

\textbf{TODO:} Report the Wilcoxon signed-rank test result here (W
statistic, $p$-value, and whether $p < 0.05$), and state plainly if the
result is not statistically significant -- this is a valid, falsifiable
outcome per the research question in Section I, not a result to hide or
obscure.

\section{Discussion}

\subsection{Failure Analysis}
\label{sec:failure-analysis}

% TODO: THIS SUBSECTION IS REQUIRED BY THE MILESTONE 4 RUBRIC AND MUST
% BE WRITTEN FROM YOUR ACTUAL CONFUSION MATRICES AND MISCLASSIFIED
% EXAMPLES. Pull 10-15 real misclassified comments (both directions:
% hate misclassified as offensive/neutral, and vice versa) and discuss
% what linguistic or contextual features likely caused the error.
% Since the primary social harm from a real deployed system comes from
% false negatives on the hate class, prioritise discussing hate-class
% recall failures specifically.

During development, a critical failure mode was identified and is worth
reporting transparently as part of this paper's methodological
contribution: an initial dataset later found to contain templated,
near-duplicate content (Section~\ref{sec:audit}) produced a macro-F1 of
1.0000 across every cross-validation fold, for both the baseline and
proposed conditions, with zero measurable difference between them. This
was diagnosed, not merely observed: a trivial bag-of-words logistic
regression model, entirely incapable of contextual language
understanding, achieved the identical ceiling score on the same data,
confirming that the near-perfect performance reflected the dataset's
narrow, near-deterministic per-class vocabulary rather than genuine
model competence. This finding is reported here as a methodological
caution for future Sinhala-English code-mixed hate speech research: near-
duplicate and template-generated content -- whether from data
augmentation or from repetitive real-world posting patterns -- can
silently invalidate cross-validation estimates unless explicitly audited
for, and macro-F1 alone, without a duplication/leakage audit, is
insufficient evidence of genuine classifier competence.

% TODO: follow the above with failure analysis of your REAL model's
% REAL errors once available.

\subsection{Ethical Reflection}
\label{sec:ethics}

This research handles sensitive content directly implicating protected
and marginalised communities in Sri Lanka, including Tamil and Muslim
minorities historically targeted during periods of communal unrest.
Several ethical safeguards were built into the methodology from the
outset. All data collection was restricted to publicly available posts
and comments; private accounts, groups, and direct messages were
excluded entirely. User identifiers -- names, profile URLs, and
usernames -- were stripped prior to annotation and are not present in
the final dataset, consistent with data minimisation principles.
Annotators were informed of the potential psychological effects of
prolonged exposure to hateful and offensive content and advised to limit
annotation sessions to 30 minutes at a time.

A specific bias risk, documented by Sap et al. \cite{sap2019}, is that
hate speech classifiers trained predominantly on majority-dialect data
can systematically misclassify minority-dialect or minority-community
speech -- for example, flagging in-group reclaimed language or
minority-community political speech as hate speech at disproportionate
rates. Clear annotation guidelines defining each class with worked
examples were used specifically to reduce annotator subjectivity and
mitigate bias against minority-community speech, though this risk cannot
be eliminated through guidelines alone and would require targeted
fairness auditing (e.g., per-community false-positive-rate analysis) as
future work before any real-world deployment.

Finally, the Online Safety Act No. 9 of 2023 (Sri Lanka) creates a
regulatory context in which automated content moderation tools for
code-mixed Sinhala-English text may see real-world deployment; the
authors caution against deploying any model resulting from this research
in a production content-moderation context without further fairness
auditing, human-in-the-loop review of flagged content, and an appeals
mechanism for affected users, given the demonstrated risk of
disproportionate impact on minority communities.

\subsection{Limitations}
\label{sec:limitations}

% TODO: state your ACTUAL final limitations honestly. Likely candidates
% based on this project's development process, to keep IF true of your
% final submission, or remove if not:
\begin{itemize}
\item Facebook data collection was constrained by Meta's post-2018
Graph API access policy, which requires Page-admin access or App Review
for comment-level data beyond a developer's own pages; % TODO: state
% whether this meant Facebook data was reduced in volume, excluded
% entirely, or obtained through an alternative access path.
\item The final annotated dataset size ( % TODO: insert real N %) is
smaller than the original target of 2,500 comments, reflecting the time
constraints of a single-semester capstone project; this limits the
statistical power of the five-fold cross-validation design, particularly
for the Wilcoxon signed-rank test (with $k=5$ paired folds, the smallest
achievable two-sided $p$-value is 0.0625, meaning even a consistent
effect across all folds may not reach conventional significance);
\item % TODO: if annotation used 2 raters with an agreement-filter
% protocol rather than the originally planned 3-annotator majority-vote
% with Cohen's Kappa >= 0.70, state this explicitly here as a scope
% reduction under time constraints.
\end{itemize}

\section{Conclusion}

% TODO: replace with your actual conclusion once results are final.
This paper presented a methodology for multi-class Sinhala-English
code-mixed hate speech detection with emoji feature integration,
addressing three gaps identified in prior work: binary-only
classification, emoji removal during preprocessing, and single-platform
data collection. % TODO: state your actual finding regarding the
% research question, and its implication for future emoji-aware,
% code-mixed hate speech detection research.

\section*{Acknowledgment}
The authors thank % TODO: module coordinator name / volunteer annotator
% name, with permission % for their contributions to this project.

\begin{thebibliography}{00}

\bibitem{muthuthanthri2024} M. Muthuthanthri and R. I. Smith, ``Hate
Speech Detection for Transliterated English and Sinhala Code-Mixed
Data,'' in \textit{Proc. ICARC}, 2024, pp. 155--160.

\bibitem{liyanage2021} O. Liyanage and K. Jayakumar, ``Hate Speech
Detection in Sinhala-English Code-Mixed Language,'' in \textit{Proc.
ICTer}, 2021, pp. 225--230.

\bibitem{hettiarachchi2020} N. Hettiarachchi, R. Weerasinghe, and R.
Pushpanda, ``Detecting Hate Speech in Romanized Sinhala,'' in
\textit{Proc. ICTer}, 2020, pp. 250--255.

\bibitem{fernando2022} W. S. S. Fernando, R. Weerasinghe, and E. R. A. D.
Bandara, ``Sinhala Hate Speech Detection Using Machine Learning and Deep
Learning,'' in \textit{Proc. ICTer}, 2022, pp. 166--171.

\bibitem{davidson2017} T. Davidson, D. Warmsley, M. Macy, and I. Weber,
``Automated Hate Speech Detection and the Problem of Offensive
Language,'' in \textit{Proc. ICWSM}, vol. 11, 2017, pp. 512--515.

\bibitem{vasantharajan2022} C. Vasantharajan and U. Thayasivam, ``Towards
Offensive Language Identification for Tamil Code-Mixed YouTube
Comments,'' \textit{SN Comput. Sci.}, vol. 3, p. 94, 2022.

\bibitem{conneau2020} A. Conneau \textit{et al.}, ``Unsupervised
Cross-Lingual Representation Learning at Scale,'' in \textit{Proc. ACL},
2020, pp. 8440--8451.

\bibitem{meedin2022} N. Meedin, M. Caldera, S. Perera, and I. Perera, ``A
Novel Annotation Scheme to Generate Hate Speech Corpus through
Crowdsourcing and Active Learning,'' \textit{Int. J. Adv. Comput. Sci.
Appl.}, vol. 13, 2022.

\bibitem{rathnayake2022} H. Rathnayake, J. Sumanapala, R. Rukshani, and
S. Ranathunga, ``Adapter-Based Fine-Tuning of Pre-Trained Multilingual
Language Models for Code-Mixed Text Classification,'' \textit{Knowl.
Inf. Syst.}, vol. 64, pp. 1937--1966, 2022.

\bibitem{munasinghe2022} S. Munasinghe and U. Thayasivam, ``A Deep
Learning Ensemble Hate Speech Detection Approach for Sinhala Tweets,''
in \textit{Proc. MERCon}, 2022, pp. 1--6.

\bibitem{mozafari2020} M. Mozafari, R. Farahbakhsh, and N. Crespi, ``A
BERT-Based Transfer Learning Approach for Hate Speech Detection in
Online Social Media,'' in \textit{Proc. Complex Networks VIII},
Springer, 2020, pp. 928--940.

\bibitem{barbieri2020} F. Barbieri, J. Camacho-Collados, L.
Espinosa-Anke, and L. Neves, ``TweetEval: Unified Benchmark and
Comparative Evaluation for Tweet Classification,'' in \textit{Proc.
EMNLP Findings}, 2020, pp. 1644--1650.

\bibitem{qureshi2021} K. A. Qureshi and M. Sabih, ``Un-Compromised
Credibility: Social Media Based Multi-Class Hate Speech Classification
for Text,'' \textit{IEEE Access}, vol. 9, pp. 109465--109477, 2021.

\bibitem{sap2019} M. Sap, D. Card, S. Gabriel, Y. Choi, and N. A. Smith,
``The Risk of Racial Bias in Hate Speech Detection,'' in \textit{Proc.
ACL}, 2019, pp. 1668--1678.

\end{thebibliography}

\end{document}
